\chapter{La méthode des panneaux de Hess-Smith}\label{ch:panneaux}
\iterbanner{V0}

\section{Le principe}
On remplace le contour par $N$ segments rectilignes, les \textbf{panneaux}. Sur chacun, on pose :
\begin{itemize}
\item une nappe de sources d'intensité \textbf{constante} $\sigma_j$, propre au panneau $j$ ;
\item une nappe tourbillonnaire d'intensité $\gamma$, constante, et \textbf{la même sur tous les panneaux}.
\end{itemize}
Cela fait $N+1$ inconnues : $\sigma_1,\dots,\sigma_N$ et $\gamma$. On écrit $N+1$ équations :
\begin{itemize}
\item la non-pénétration $\vect V\cdot\vect n_i=0$ au milieu de chaque panneau, appelé \textbf{point de collocation} : $N$ équations ;
\item la condition de Kutta : 1 équation.
\end{itemize}

\paragraph{Pourquoi des sources et un seul tourbillon ?} Les sources seules ont une circulation totale nulle, donc ne portent pas (Kutta-Joukowski). Il faut de la circulation, mais un tourbillon libre par panneau donnerait $2N$ inconnues pour $N+1$ équations : le problème serait sous-déterminé. Hess et Smith (1967) prennent le minimum : une seule intensité tourbillonnaire uniforme, que la condition de Kutta fixe. C'est la formulation la plus simple qui donne une solution unique avec portance.

\section{Géométrie des panneaux}
Le panneau $j$ va du point $(x_j,y_j)$ au point $(x_{j+1},y_{j+1})$ du contour. On en tire (figure~\ref{fig:panneau}) :
\[
L_j=\sqrt{(x_{j+1}-x_j)^2+(y_{j+1}-y_j)^2},\qquad
\vect t_j=\frac{(x_{j+1}-x_j,\ y_{j+1}-y_j)}{L_j},\qquad
\vect n_j=(t_{y,j},\,-t_{x,j}),
\]
et le point de collocation $(\bar x_j,\bar y_j)$ au milieu. Comme le contour est parcouru dans le sens trigonométrique, $\vect n_j$ (à droite de $\vect t_j$) pointe vers l'extérieur, c'est-à-dire vers le fluide.

\begin{figure}[htbp]
\centering
\begin{tikzpicture}[scale=1.1]
% panneau
\draw[cink,very thick] (0,0) -- (4,0);
\fill[cink] (0,0) circle (1.6pt) node[below left,font=\small] {début $(x_j,y_j)$};
\fill[cink] (4,0) circle (1.6pt) node[below right,font=\small] {fin $(x_{j+1},y_{j+1})$};
\fill[cgreen] (2,0) circle (1.8pt) node[below left=3pt and 1pt,font=\small,cgreen] {collocation};
\draw[arr,cblue] (2,0) -- (3,0) node[above,font=\small,cblue] {$\vect t_j$};
\draw[arr,cblue] (2,0) -- (2,-0.9) node[right,font=\small,cblue] {$\vect n_j$ (vers le fluide)};
\draw[arr,cmuted] (2,0) -- (2,0.7) node[left,font=\small,cmuted] {$\vect e_\eta$};
% repère local
\draw[arr,cmuted] (0,0) -- (0.9,0) ; \node[cmuted,font=\small] at (0.8,0.2) {$\xi$};
% point P
\coordinate (P) at (2.9,2.2);
\fill[corange] (P) circle (1.8pt) node[above right,font=\small,corange] {$P=(\xi,\eta)$};
\draw[corange,thin] (0,0) -- (P) node[pos=0.5,above left,font=\small] {$r_1$};
\draw[corange,thin] (4,0) -- (P) node[pos=0.5,right,font=\small] {$r_2$};
\draw[corange] (0.7,0) arc (0:37.2:0.7); \node[corange,font=\small] at (0.95,0.32) {$\theta_1$};
\draw[corange] (4.5,0) arc (0:116.6:0.5); \node[corange,font=\small] at (4.55,0.62) {$\theta_2$};
\draw[cmuted,dashed] (4,0) -- (4.7,0);
\node[font=\small,cmuted,align=left] at (-1.5,1.3) {intérieur du profil\\ (côté $\eta>0$)};
\node[font=\small,cmuted] at (-1.2,-1.0) {fluide ($\eta<0$)};
\end{tikzpicture}
\caption{Un panneau et son repère local. $\xi$ suit la tangente, $\eta$ la normale gauche, qui pointe vers l'intérieur du profil. Les angles $\theta_1$, $\theta_2$ sont ceux sous lesquels les extrémités voient le point $P$.}\label{fig:panneau}
\end{figure}

\section{La vitesse induite par un panneau}
Plaçons le panneau sur l'axe $\xi\in[0,L]$, $\eta=0$, et calculons la vitesse qu'il induit en un point $P=(\xi,\eta)$.

\paragraph{Source.} Un élément $\dd\xi'$ en $(\xi',0)$ est une source de débit $\sigma\,\dd\xi'$. Sa vitesse en $P$ est radiale :
\[
\dd\vect V=\frac{\sigma\,\dd\xi'}{2\pi}\,\frac{(\xi-\xi',\ \eta)}{(\xi-\xi')^2+\eta^2}.
\]
En intégrant :
\[
u_\sigma=\frac{\sigma}{2\pi}\int_0^L\frac{\xi-\xi'}{(\xi-\xi')^2+\eta^2}\,\dd\xi'
=\frac{\sigma}{4\pi}\Big[-\ln\big((\xi-\xi')^2+\eta^2\big)\Big]_0^L
=\frac{\sigma}{2\pi}\ln\frac{r_1}{r_2},
\]
\[
v_\sigma=\frac{\sigma}{2\pi}\int_0^L\frac{\eta}{(\xi-\xi')^2+\eta^2}\,\dd\xi'=\frac{\sigma}{2\pi}\big(\theta_2-\theta_1\big),
\qquad \theta_k=\operatorname{atan2}(\eta,\ \xi-\xi_k).
\]
La composante tangentielle dépend du rapport des distances aux extrémités, la composante normale de l'angle sous lequel on voit le panneau.

\paragraph{Tourbillon.} La vitesse d'un tourbillon est celle d'une source tournée de $+90^\circ$ : $(u,v)\to(-v,u)$. On obtient sans nouveau calcul
\[
u_\gamma=-\frac{\gamma}{2\pi}\big(\theta_2-\theta_1\big),\qquad v_\gamma=\frac{\gamma}{2\pi}\ln\frac{r_1}{r_2}.
\]

\paragraph{Auto-influence.} Quand $P$ est le point de collocation du panneau lui-même, $r_1=r_2$ donc $\ln(r_1/r_2)=0$. Pour l'angle, il faut savoir de quel côté on s'approche. Le fluide est du côté $\eta<0$. En faisant $\eta\to0^-$ avec $0<\xi<L$, on a $\theta_1\to0$ et $\theta_2\to-\pi$, donc $\theta_2-\theta_1=-\pi$. On retrouve les demi-sauts du chapitre précédent :
\begin{itemize}
\item la source induit sur elle-même $v_\sigma=-\sigma/2$ selon $\vect e_\eta$, soit $+\sigma/2$ selon la normale sortante : une source positive « souffle » vers le fluide ;
\item le tourbillon induit $u_\gamma=+\gamma/2$ selon la tangente.
\end{itemize}

\paragraph{Retour au repère global.} Avec $\vect e_\xi=\vect t_j=(t_x,t_y)$ et $\vect e_\eta=(-t_y,t_x)$, une vitesse locale $(u,v)$ vaut $(u\,t_x-v\,t_y,\ u\,t_y+v\,t_x)$ dans le repère global.

\begin{codebox}[ : \code{panel.influence\_matrices}]
Tout est vectorisé : on calcule d'un coup les $N\times N$ couples (collocation $i$, panneau $j$).
\begin{lstlisting}
dx = P.xc[:, None] - P.xs[None, :]          # collocation i - debut du panneau j
dy = P.yc[:, None] - P.ys[None, :]
xi  =  dx * tx + dy * ty                    # coordonnees locales du panneau j
eta = -dx * ty + dy * tx
r1, r2 = np.hypot(xi, eta), np.hypot(xi - L, eta)
lnr = np.log(r1 / r2)
dth = np.arctan2(eta, xi - L) - np.arctan2(eta, xi)    # theta2 - theta1
idx = np.arange(P.n)
lnr[idx, idx] = 0.0                         # auto-influence, cote fluide
dth[idx, idx] = -np.pi
us_l, vs_l = lnr / (2*np.pi), dth / (2*np.pi)      # source unitaire
uv_l, vv_l = -dth / (2*np.pi), lnr / (2*np.pi)     # tourbillon unitaire
\end{lstlisting}
Les quatre matrices renvoyées (\code{usx, usy, uvx, uvy}) contiennent les composantes globales de la vitesse induite en $i$ par une source unitaire et par un tourbillon unitaire posés sur le panneau $j$. Elles ne dépendent que de la géométrie, pas de l'incidence.
\end{codebox}

\section{Le système linéaire}
Notons $A^n_{ij}$ et $A^t_{ij}$ les composantes normale ($\cdot\,\vect n_i$) et tangentielle ($\cdot\,\vect t_i$) de la vitesse induite en $i$ par une source unitaire sur le panneau $j$, et $B^n_{ij}$, $B^t_{ij}$ les mêmes pour un tourbillon unitaire. La vitesse au point de collocation $i$ s'écrit
\[
\vect V_i=\vect U_\infty+\sum_j\sigma_j\,(\text{vitesse de la source }j)+\gamma\sum_j(\text{vitesse du tourbillon }j),
\qquad \vect U_\infty=(\cos\alpha,\ \sin\alpha).
\]

\paragraph{Non-pénétration.} Pour $i=1,\dots,N$ :
\[
\sum_{j=1}^N A^n_{ij}\,\sigma_j+\Big(\sum_{j=1}^N B^n_{ij}\Big)\gamma=-\vect U_\infty\cdot\vect n_i .
\]

\paragraph{Kutta discrète.} Les deux panneaux qui touchent le bord de fuite sont le premier (extrados, $\vect t_1$ dirigé vers l'amont) et le dernier (intrados, $\vect t_N$ dirigé vers l'aval). Leurs tangentes sont donc presque opposées. L'écoulement réel quitte le bord de fuite vers l'aval des deux côtés : $V_{t,1}<0$ et $V_{t,N}>0$. Demander des vitesses de même norme revient à
\begin{gather*}
V_{t,1}+V_{t,N}=0\\
\Longleftrightarrow\quad
\sum_j\big(A^t_{1j}+A^t_{Nj}\big)\sigma_j+\Big(\sum_jB^t_{1j}+\sum_jB^t_{Nj}\Big)\gamma=-\vect U_\infty\cdot(\vect t_1+\vect t_N).
\end{gather*}

\paragraph{Forme matricielle.}
\[
\begin{pmatrix}
A^n & B^n\vect 1\\[2pt]
A^t_{1\cdot}+A^t_{N\cdot} & \vect 1^\top\big(B^t_{1\cdot}+B^t_{N\cdot}\big)^\top
\end{pmatrix}
\begin{pmatrix}\boldsymbol\sigma\\ \gamma\end{pmatrix}
=-\begin{pmatrix}\big(\vect U_\infty\cdot\vect n_i\big)_{i}\\ \vect U_\infty\cdot(\vect t_1+\vect t_N)\end{pmatrix}.
\]
C'est un système \textbf{dense} de taille $N+1=201$. La matrice ne dépend que de la géométrie ; seul le second membre dépend de $\alpha$.

\begin{piege}
Écrire $V_{t,1}=V_{t,N}$ au lieu de $V_{t,1}+V_{t,N}=0$ est l'erreur classique : avec des tangentes opposées, cela imposerait des vitesses physiques \emph{opposées} au bord de fuite. Le symptôme est un $C_L$ faux et un $C_{L,C_p}$ très différent de $C_{L,KJ}$.
\end{piege}

\begin{codebox}[ : \code{panel.solve}]
\begin{lstlisting}
A = np.zeros((N + 1, N + 1))
A[:N, :N] = Asn                          # sources -> vitesse normale
A[:N, N]  = Avn.sum(axis=1)              # tourbillon uniforme -> vitesse normale
A[N, :N]  = Ast[0] + Ast[-1]             # Kutta : Vt_1 + Vt_N = 0
A[N, N]   = Avt[0].sum() + Avt[-1].sum()
lu = np.linalg.inv(A)                    # reutilisable pour toute incidence
b[:N] = -(Uinf[0] * P.nx + Uinf[1] * P.ny)
b[N]  = -(Uinf[0] * (P.tx[0] + P.tx[-1]) + Uinf[1] * (P.ty[0] + P.ty[-1]))
q = lu @ b;  sigma, gamma = q[:N], q[N]
\end{lstlisting}
Le dictionnaire optionnel \code{cache} conserve les panneaux, les matrices d'influence et l'inverse. Changer d'incidence ne coûte alors qu'un produit matrice-vecteur, ce que \code{alpha\_for\_cl} exploite (ch.~\ref{ch:modele}). L'inverse explicite est acceptable ici : $201\times201$, une fois par profil.
\end{codebox}

\section{Ce qu'on tire de la solution}
\paragraph{Vitesse tangentielle et pression.} Aux points de collocation,
\[
V_{t,i}=\vect U_\infty\cdot\vect t_i+\sum_jA^t_{ij}\sigma_j+\gamma\sum_jB^t_{ij},\qquad C_{p,i}=1-V_{t,i}^2 .
\]
$V_{t,i}$ est \emph{signée} selon $\vect t_i$ : elle est négative sur l'extrados (l'écoulement va du bord d'attaque vers le bord de fuite, à l'opposé du sens de parcours) et positive sur l'intrados. Le changement de signe repère le point d'arrêt (ch.~\ref{ch:cl}).

\paragraph{Efforts.} La force de pression sur le panneau $i$ vaut $-C_{p,i}\,\vect n_i\,L_i$ (en unités de $\tfrac12\rho U_\infty^2c$). On somme, puis on projette :
\[
\vect F=-\sum_iC_{p,i}\,\vect n_i\,L_i,\qquad
C_L=F_y\cos\alpha-F_x\sin\alpha,\qquad C_{D,p}=F_x\cos\alpha+F_y\sin\alpha,
\]
\[
C_{M,c/4}=-\sum_i\Big[(\bar x_i-\tfrac14)\,f_{y,i}-\bar y_i\,f_{x,i}\Big]\qquad\text{(positif à cabrer).}
\]

\paragraph{Deux contrôles gratuits.} La circulation totale de la nappe tourbillonnaire vaut $\gamma$ fois le périmètre, comptée dans le sens trigonométrique. La portance correspond à une circulation horaire $\Gamma=-\gamma\,\ell_{tot}$, d'où $C_{L,KJ}=-2\gamma\,\ell_{tot}$ par Kutta-Joukowski. On doit avoir $C_{L,KJ}\approx C_{L,C_p}$, et $C_{D,p}\to0$ (d'Alembert). Les deux écarts mesurent l'erreur de discrétisation.

\begin{figure}[htbp]
\centering
% figure générée par tools/make_figures.py : cp
\begin{tikzpicture}
\begin{groupplot}[group style={group size=3 by 1,horizontal sep=0.9cm,y descriptions at=edge left},
  width=0.34\linewidth,height=5.2cm,y dir=reverse,xmin=0,xmax=1,ymin=-2.6,ymax=1.1,
  xlabel={$x/c$},ylabel={$C_p$},title style={font=\small},tick label style={font=\scriptsize}]
\nextgroupplot[title={NACA 0012, $\alpha=0^\circ$}]
\addplot[cblue,very thick] coordinates {(0.00012336,0.98707) (0.00061668,0.88936) (0.0016028,0.72336) (0.0030808,0.53089) (0.0050492,0.34529) (0.0075061,0.18392) (0.010449,0.051559) (0.013875,-0.05372) (0.017781,-0.13625) (0.022162,-0.2006) (0.027016,-0.25076) (0.032336,-0.28994) (0.038117,-0.3206) (0.044355,-0.34463) (0.051042,-0.36343) (0.058172,-0.37807) (0.065738,-0.38937) (0.073733,-0.39793) (0.082148,-0.40422) (0.090976,-0.40862) (0.10021,-0.4114) (0.10983,-0.4128) (0.11984,-0.41299) (0.13023,-0.41212) (0.14098,-0.4103) (0.15209,-0.40765) (0.16354,-0.40424) (0.17532,-0.40015) (0.18742,-0.39545) (0.19983,-0.39018) (0.21253,-0.38441) (0.22552,-0.37819) (0.23878,-0.37155) (0.2523,-0.36454) (0.26606,-0.35719) (0.28006,-0.34955) (0.29427,-0.34165) (0.30868,-0.33352) (0.32328,-0.32519) (0.33806,-0.3167) (0.353,-0.30807) (0.36808,-0.29932) (0.38329,-0.29048) (0.39862,-0.28157) (0.41405,-0.27261) (0.42956,-0.26363) (0.44514,-0.25463) (0.46078,-0.24563) (0.47645,-0.23665) (0.49215,-0.2277) (0.50785,-0.21878) (0.52355,-0.2099) (0.53922,-0.20107) (0.55486,-0.19229) (0.57044,-0.18356) (0.58595,-0.17489) (0.60138,-0.16627) (0.61671,-0.15769) (0.63192,-0.14916) (0.647,-0.14067) (0.66194,-0.13221) (0.67672,-0.12377) (0.69132,-0.11535) (0.70573,-0.10693) (0.71994,-0.098492) (0.73394,-0.090035) (0.7477,-0.081542) (0.76122,-0.072998) (0.77448,-0.064388) (0.78747,-0.055696) (0.80017,-0.046906) (0.81258,-0.037999) (0.82468,-0.028959) (0.83646,-0.019766) (0.84791,-0.010401) (0.85902,-0.00084368) (0.86977,0.0089278) (0.88016,0.018936) (0.89017,0.029206) (0.89979,0.039764) (0.90902,0.050638) (0.91785,0.061862) (0.92627,0.073471) (0.93426,0.085506) (0.94183,0.098016) (0.94896,0.11106) (0.95565,0.12469) (0.96188,0.13901) (0.96766,0.15411) (0.97298,0.17012) (0.97784,0.18719) (0.98222,0.20556) (0.98612,0.2255) (0.98955,0.24743) (0.99249,0.27195) (0.99495,0.29995) (0.99692,0.33298) (0.9984,0.37378) (0.99938,0.42821) (0.99988,0.50906)};
\addplot[corange,very thick] coordinates {(0.00012336,0.98707) (0.00061668,0.88936) (0.0016028,0.72336) (0.0030808,0.53089) (0.0050492,0.34529) (0.0075061,0.18392) (0.010449,0.051559) (0.013875,-0.05372) (0.017781,-0.13625) (0.022162,-0.2006) (0.027016,-0.25076) (0.032336,-0.28994) (0.038117,-0.3206) (0.044355,-0.34463) (0.051042,-0.36343) (0.058172,-0.37807) (0.065738,-0.38937) (0.073733,-0.39793) (0.082148,-0.40422) (0.090976,-0.40862) (0.10021,-0.4114) (0.10983,-0.4128) (0.11984,-0.41299) (0.13023,-0.41212) (0.14098,-0.4103) (0.15209,-0.40765) (0.16354,-0.40424) (0.17532,-0.40015) (0.18742,-0.39545) (0.19983,-0.39018) (0.21253,-0.38441) (0.22552,-0.37819) (0.23878,-0.37155) (0.2523,-0.36454) (0.26606,-0.35719) (0.28006,-0.34955) (0.29427,-0.34165) (0.30868,-0.33352) (0.32328,-0.32519) (0.33806,-0.3167) (0.353,-0.30807) (0.36808,-0.29932) (0.38329,-0.29048) (0.39862,-0.28157) (0.41405,-0.27261) (0.42956,-0.26363) (0.44514,-0.25463) (0.46078,-0.24563) (0.47645,-0.23665) (0.49215,-0.2277) (0.50785,-0.21878) (0.52355,-0.2099) (0.53922,-0.20107) (0.55486,-0.19229) (0.57044,-0.18356) (0.58595,-0.17489) (0.60138,-0.16627) (0.61671,-0.15769) (0.63192,-0.14916) (0.647,-0.14067) (0.66194,-0.13221) (0.67672,-0.12377) (0.69132,-0.11535) (0.70573,-0.10693) (0.71994,-0.098492) (0.73394,-0.090035) (0.7477,-0.081542) (0.76122,-0.072998) (0.77448,-0.064388) (0.78747,-0.055696) (0.80017,-0.046906) (0.81258,-0.037999) (0.82468,-0.028959) (0.83646,-0.019766) (0.84791,-0.010401) (0.85902,-0.00084368) (0.86977,0.0089278) (0.88016,0.018936) (0.89017,0.029206) (0.89979,0.039764) (0.90902,0.050638) (0.91785,0.061862) (0.92627,0.073471) (0.93426,0.085506) (0.94183,0.098016) (0.94896,0.11106) (0.95565,0.12469) (0.96188,0.13901) (0.96766,0.15411) (0.97298,0.17012) (0.97784,0.18719) (0.98222,0.20556) (0.98612,0.2255) (0.98955,0.24743) (0.99249,0.27195) (0.99495,0.29995) (0.99692,0.33298) (0.9984,0.37378) (0.99938,0.42821) (0.99988,0.50906)};
\node[font=\scriptsize,anchor=south east,fill=white,inner sep=1pt] at (rel axis cs:0.98,0.02) {$C_L=0.000$};
\nextgroupplot[title={NACA 0012, $\alpha=5^\circ$}]
\addplot[cblue,very thick] coordinates {(0.00012336,-0.64686) (0.00061668,-1.1651) (0.0016028,-1.5826) (0.0030808,-1.8604) (0.0050492,-2.0079) (0.0075061,-2.058) (0.010449,-2.0448) (0.013875,-1.9949) (0.017781,-1.9261) (0.022162,-1.8492) (0.027016,-1.7707) (0.032336,-1.694) (0.038117,-1.6208) (0.044355,-1.5519) (0.051042,-1.4874) (0.058172,-1.4272) (0.065738,-1.3712) (0.073733,-1.3188) (0.082148,-1.2698) (0.090976,-1.2239) (0.10021,-1.1806) (0.10983,-1.1397) (0.11984,-1.101) (0.13023,-1.0642) (0.14098,-1.0291) (0.15209,-0.99554) (0.16354,-0.9633) (0.17532,-0.93229) (0.18742,-0.90238) (0.19983,-0.87348) (0.21253,-0.84548) (0.22552,-0.81832) (0.23878,-0.79194) (0.2523,-0.76627) (0.26606,-0.74128) (0.28006,-0.71692) (0.29427,-0.69316) (0.30868,-0.66997) (0.32328,-0.64733) (0.33806,-0.62521) (0.353,-0.60361) (0.36808,-0.5825) (0.38329,-0.56186) (0.39862,-0.54169) (0.41405,-0.52198) (0.42956,-0.50271) (0.44514,-0.48387) (0.46078,-0.46545) (0.47645,-0.44744) (0.49215,-0.42982) (0.50785,-0.41259) (0.52355,-0.39574) (0.53922,-0.37924) (0.55486,-0.36309) (0.57044,-0.34726) (0.58595,-0.33176) (0.60138,-0.31655) (0.61671,-0.30162) (0.63192,-0.28696) (0.647,-0.27254) (0.66194,-0.25836) (0.67672,-0.24438) (0.69132,-0.23058) (0.70573,-0.21696) (0.71994,-0.20348) (0.73394,-0.19013) (0.7477,-0.17688) (0.76122,-0.16372) (0.77448,-0.15062) (0.78747,-0.13756) (0.80017,-0.12451) (0.81258,-0.11146) (0.82468,-0.098382) (0.83646,-0.085252) (0.84791,-0.072049) (0.85902,-0.058747) (0.86977,-0.045323) (0.88016,-0.031751) (0.89017,-0.018004) (0.89979,-0.0040549) (0.90902,0.010128) (0.91785,0.024578) (0.92627,0.039331) (0.93426,0.05443) (0.94183,0.069922) (0.94896,0.085863) (0.95565,0.10232) (0.96188,0.11937) (0.96766,0.13712) (0.97298,0.15569) (0.97784,0.17524) (0.98222,0.19598) (0.98612,0.21819) (0.98955,0.24228) (0.99249,0.26883) (0.99495,0.29873) (0.99692,0.33345) (0.9984,0.37572) (0.99938,0.4313) (0.99988,0.51279)};
\addplot[corange,very thick] coordinates {(0.00012336,-0.11677) (0.00061668,0.34601) (0.0016028,0.68737) (0.0030808,0.8933) (0.0050492,0.98506) (0.0075061,0.99738) (0.010449,0.96181) (0.013875,0.90101) (0.017781,0.82926) (0.022162,0.75485) (0.027016,0.68226) (0.032336,0.61371) (0.038117,0.55015) (0.044355,0.49181) (0.051042,0.43857) (0.058172,0.39015) (0.065738,0.34616) (0.073733,0.30624) (0.082148,0.27002) (0.090976,0.23715) (0.10021,0.20732) (0.10983,0.18026) (0.11984,0.15573) (0.13023,0.1335) (0.14098,0.11337) (0.15209,0.095174) (0.16354,0.078756) (0.17532,0.063974) (0.18742,0.050701) (0.19983,0.038822) (0.21253,0.028232) (0.22552,0.018834) (0.23878,0.010539) (0.2523,0.0032645) (0.26606,-0.0030652) (0.28006,-0.0085219) (0.29427,-0.013172) (0.30868,-0.017079) (0.32328,-0.020299) (0.33806,-0.022889) (0.353,-0.024899) (0.36808,-0.026378) (0.38329,-0.027369) (0.39862,-0.027916) (0.41405,-0.028056) (0.42956,-0.027826) (0.44514,-0.02726) (0.46078,-0.026387) (0.47645,-0.025235) (0.49215,-0.02383) (0.50785,-0.022195) (0.52355,-0.020349) (0.53922,-0.018311) (0.55486,-0.016096) (0.57044,-0.013718) (0.58595,-0.011188) (0.60138,-0.008514) (0.61671,-0.0057039) (0.63192,-0.0027626) (0.647,0.00030676) (0.66194,0.0035028) (0.67672,0.0068259) (0.69132,0.010278) (0.70573,0.013862) (0.71994,0.017583) (0.73394,0.021448) (0.7477,0.025463) (0.76122,0.029637) (0.77448,0.033981) (0.78747,0.038505) (0.80017,0.04322) (0.81258,0.048141) (0.82468,0.053281) (0.83646,0.058655) (0.84791,0.06428) (0.85902,0.070173) (0.86977,0.076354) (0.88016,0.082842) (0.89017,0.089661) (0.89979,0.096835) (0.90902,0.10439) (0.91785,0.11236) (0.92627,0.12078) (0.93426,0.1297) (0.94183,0.13914) (0.94896,0.14919) (0.95565,0.15989) (0.96188,0.17134) (0.96766,0.18363) (0.97298,0.19689) (0.97784,0.2113) (0.98222,0.22706) (0.98612,0.24447) (0.98955,0.26394) (0.99249,0.28607) (0.99495,0.31179) (0.99692,0.34262) (0.9984,0.38134) (0.99938,0.4338) (0.99988,0.51279)};
\node[font=\scriptsize,anchor=south east,fill=white,inner sep=1pt] at (rel axis cs:0.98,0.02) {$C_L=0.603$};
\nextgroupplot[title={NACA 2412, $\alpha=2^\circ$}]
\addplot[cblue,very thick] coordinates {(-1.4838e-05,0.77064) (0.00020448,0.52203) (0.00092198,0.24185) (0.0021378,-0.021587) (0.0038514,-0.24115) (0.0060621,-0.41069) (0.0087685,-0.53586) (0.011969,-0.62605) (0.01566,-0.69031) (0.01984,-0.73587) (0.024505,-0.76813) (0.029651,-0.79094) (0.035273,-0.807) (0.041368,-0.81819) (0.047928,-0.82581) (0.054949,-0.83077) (0.062423,-0.8337) (0.070344,-0.83503) (0.078703,-0.83507) (0.087494,-0.83403) (0.096707,-0.83205) (0.10633,-0.82924) (0.11636,-0.82565) (0.12679,-0.82133) (0.13759,-0.81631) (0.14877,-0.81062) (0.16031,-0.80426) (0.17219,-0.79724) (0.18442,-0.78957) (0.19696,-0.78125) (0.20982,-0.77229) (0.22297,-0.76269) (0.2364,-0.75246) (0.25011,-0.74159) (0.26407,-0.73009) (0.27826,-0.71795) (0.29268,-0.70518) (0.30731,-0.69175) (0.32214,-0.67766) (0.33714,-0.66285) (0.3523,-0.64726) (0.36761,-0.63075) (0.38305,-0.61298) (0.39857,-0.59301) (0.41414,-0.57296) (0.42974,-0.55449) (0.44542,-0.53715) (0.46115,-0.52047) (0.47691,-0.50423) (0.49269,-0.48834) (0.50848,-0.47272) (0.52425,-0.45732) (0.54,-0.44213) (0.5557,-0.4271) (0.57135,-0.41223) (0.58692,-0.39749) (0.6024,-0.38287) (0.61777,-0.36835) (0.63302,-0.35393) (0.64814,-0.33957) (0.6631,-0.32528) (0.6779,-0.31104) (0.69252,-0.29683) (0.70695,-0.28263) (0.72116,-0.26844) (0.73515,-0.25424) (0.74891,-0.24002) (0.76242,-0.22575) (0.77566,-0.21142) (0.78863,-0.19702) (0.80131,-0.18253) (0.81369,-0.16793) (0.82576,-0.15322) (0.8375,-0.13837) (0.84891,-0.12336) (0.85997,-0.10818) (0.87068,-0.092813) (0.88102,-0.077237) (0.89098,-0.061432) (0.90055,-0.045378) (0.90973,-0.029051) (0.9185,-0.012424) (0.92686,0.0045331) (0.9348,0.021855) (0.94232,0.039582) (0.9494,0.057764) (0.95603,0.07646) (0.96222,0.095743) (0.96795,0.1157) (0.97323,0.13646) (0.97804,0.15817) (0.98238,0.18102) (0.98625,0.20531) (0.98965,0.2314) (0.99257,0.25988) (0.995,0.29161) (0.99695,0.32805) (0.99841,0.37189) (0.99939,0.42886) (0.99988,0.51151)};
\addplot[corange,very thick] coordinates {(0.00026156,0.9374) (0.0010289,0.9995) (0.0022837,0.96529) (0.0040239,0.86694) (0.006247,0.74009) (0.0089501,0.6104) (0.01213,0.49149) (0.015782,0.3885) (0.019902,0.30194) (0.024485,0.23031) (0.029527,0.17147) (0.035021,0.1233) (0.040961,0.083896) (0.047342,0.051689) (0.054155,0.025391) (0.061395,0.0039608) (0.069053,-0.013434) (0.077121,-0.027462) (0.085592,-0.038666) (0.094457,-0.047483) (0.10371,-0.054271) (0.11333,-0.059326) (0.12333,-0.062893) (0.13367,-0.065179) (0.14437,-0.066358) (0.1554,-0.066581) (0.16676,-0.065978) (0.17844,-0.064664) (0.19042,-0.062738) (0.20269,-0.060293) (0.21525,-0.05741) (0.22808,-0.054165) (0.24116,-0.050628) (0.25449,-0.046865) (0.26806,-0.042938) (0.28185,-0.038909) (0.29585,-0.034835) (0.31005,-0.03078) (0.32443,-0.026808) (0.33899,-0.022992) (0.3537,-0.019422) (0.36855,-0.016223) (0.38354,-0.013619) (0.39866,-0.012141) (0.41396,-0.011327) (0.42937,-0.009486) (0.44486,-0.0067567) (0.4604,-0.0035765) (0.47599,-0.00011111) (0.4916,0.0035457) (0.50723,0.0073323) (0.52285,0.011205) (0.53845,0.015132) (0.55402,0.01909) (0.56954,0.023059) (0.58499,0.027027) (0.60037,0.030985) (0.61565,0.034925) (0.63082,0.038845) (0.64587,0.042742) (0.66078,0.04662) (0.67553,0.05048) (0.69012,0.054327) (0.70452,0.058169) (0.71872,0.062013) (0.73272,0.065869) (0.74649,0.069748) (0.76001,0.073663) (0.77329,0.077626) (0.7863,0.081654) (0.79903,0.08576) (0.81147,0.089963) (0.82361,0.09428) (0.83542,0.09873) (0.84691,0.10333) (0.85806,0.10812) (0.86886,0.11309) (0.8793,0.1183) (0.88936,0.12375) (0.89904,0.12949) (0.90832,0.13554) (0.9172,0.14194) (0.92567,0.14873) (0.93372,0.15596) (0.94134,0.16368) (0.94852,0.17195) (0.95526,0.18085) (0.96155,0.19046) (0.96737,0.2009) (0.97274,0.2123) (0.97763,0.22482) (0.98205,0.23871) (0.986,0.25425) (0.98945,0.27189) (0.99242,0.29222) (0.9949,0.3162) (0.99689,0.34537) (0.99838,0.38256) (0.99938,0.43365) (0.99988,0.51151)};
\node[font=\scriptsize,anchor=south east,fill=white,inner sep=1pt] at (rel axis cs:0.98,0.02) {$C_L=0.498$};
\end{groupplot}
\end{tikzpicture}

\caption{Distributions de $C_p$ calculées par \code{panel.solve} (200 panneaux). Bleu : extrados ; orange : intrados. Axe vertical inversé, convention usuelle. Le pic d'aspiration du NACA 0012 passe de $C_p=-0{,}41$ à $\alpha=0$ à $-2{,}06$ à $\alpha=5^\circ$.}\label{fig:cp}
\end{figure}

\section{Vérifications}
\paragraph{Convergence.} La figure~\ref{fig:convergence} montre, pour le NACA 0012 à $5^\circ$, que les trois indicateurs décroissent à peu près en $1/N$ (pentes mesurées $-1{,}17$ pour $C_L$ et $-1{,}24$ pour $C_{D,p}$). C'est l'ordre attendu pour des panneaux plans à singularités constantes avec collocation au milieu. À 200 panneaux, $C_L$ est à $0{,}06$~\% de sa valeur limite.

\begin{figure}[htbp]
\centering
% figure générée par tools/make_figures.py : convergence
\begin{tikzpicture}
\begin{loglogaxis}[width=0.72\linewidth,height=6cm,xlabel={nombre de panneaux $N$},ylabel={erreur},
  legend style={draw=none,font=\small,at={(0.02,0.02)},anchor=south west},xmin=15,xmax=800,xtick={20,50,100,200,500},xticklabels={20,50,100,200,500}]
\addplot[cblue,very thick,mark=*,mark size=1.5pt] coordinates {(20,0.001539) (30,0.003499) (40,0.003447) (60,0.002396) (80,0.001768) (120,0.001137) (160,0.0008237) (200,0.000638) (320,0.0003657) (480,0.0002197) (640,0.0001488)};
\addlegendentry{$|C_L(N)-C_L^\infty|/C_L^\infty$}
\addplot[corange,very thick,mark=square*,mark size=1.5pt] coordinates {(20,0.01875) (30,0.01904) (40,0.01595) (60,0.01136) (80,0.008758) (120,0.006) (160,0.004562) (200,0.00368) (320,0.002328) (480,0.001562) (640,0.001176)};
\addlegendentry{$|C_{L,C_p}-C_{L,KJ}|/C_L^\infty$}
\addplot[cgreen,very thick,mark=triangle*,mark size=1.8pt] coordinates {(20,0.005424) (30,0.002714) (40,0.001664) (60,0.0008615) (80,0.0005582) (120,0.0003175) (160,0.0002189) (200,0.0001662) (320,9.574e-05) (480,6.09e-05) (640,4.459e-05)};
\addlegendentry{$|C_{D,p}|$ (d'Alembert)}
\addplot[cmuted,dashed,domain=40:640] {0.6/x};
\node[cmuted,font=\scriptsize] at (axis cs:330,0.0035) {pente $-1$};
\end{loglogaxis}
\end{tikzpicture}

\caption{Convergence en maillage sur le NACA 0012 à $\alpha=5^\circ$. La référence $C_L^\infty$ est calculée avec 2\,400 panneaux.}\label{fig:convergence}
\end{figure}

\paragraph{Pente et angle de portance nulle.} La figure~\ref{fig:cl-alpha} compare les panneaux à la théorie des profils minces. Les pentes mesurées valent $6{,}88$ à $6{,}89$~/rad, contre $2\pi(1+0{,}77\times0{,}12)=6{,}86$ attendu avec la correction d'épaisseur. Les angles de portance nulle valent $-2{,}13^\circ$ (2412) et $-4{,}26^\circ$ (4412), contre $-2{,}08^\circ$ et $-4{,}15^\circ$ en profil mince.

\begin{figure}[htbp]
\centering
% figure générée par tools/make_figures.py : cl_alpha
\begin{tikzpicture}
\begin{axis}[width=0.75\linewidth,height=6.4cm,xlabel={$\alpha$ (°)},ylabel={$C_L$ (non visqueux)},
  xmin=-4,xmax=10,legend columns=2,legend style={draw=none,font=\scriptsize,at={(0.5,-0.22)},anchor=north},legend cell align=left]
\addplot[cmuted,very thick,mark=*,mark size=1.3pt] coordinates {(-4,-0.4823) (-3,-0.36185) (-2,-0.2413) (-1,-0.12067) (0,5.5511e-17) (1,0.12067) (2,0.2413) (3,0.36185) (4,0.4823) (5,0.6026) (6,0.72271) (7,0.84261) (8,0.96224) (9,1.0816) (10,1.2006)};\addlegendentry{NACA 0012 (panneaux)}
\addplot[cmuted,dashed,domain=-4:10] {2*pi*(x-(-0.0000))*pi/180};\addlegendentry{profil mince, $\alpha_{L0}=-0.00^\circ$}
\addplot[cblue,very thick,mark=*,mark size=1.3pt] coordinates {(-4,-0.2265) (-3,-0.10581) (-2,0.014921) (-1,0.13565) (0,0.25635) (1,0.37698) (2,0.4975) (3,0.61788) (4,0.73808) (5,0.85806) (6,0.97778) (7,1.0972) (8,1.2163) (9,1.3351) (10,1.4534)};\addlegendentry{NACA 2412 (panneaux)}
\addplot[cblue,dashed,domain=-4:10] {2*pi*(x-(-2.0772))*pi/180};\addlegendentry{profil mince, $\alpha_{L0}=-2.08^\circ$}
\addplot[corange,very thick,mark=*,mark size=1.3pt] coordinates {(-4,0.028327) (-3,0.14925) (-2,0.27014) (-1,0.39097) (0,0.51169) (1,0.63227) (2,0.75268) (3,0.87287) (4,0.99282) (5,1.1125) (6,1.2318) (7,1.3508) (8,1.4694) (9,1.5875) (10,1.7052)};\addlegendentry{NACA 4412 (panneaux)}
\addplot[corange,dashed,domain=-4:10] {2*pi*(x-(-4.1545))*pi/180};\addlegendentry{profil mince, $\alpha_{L0}=-4.15^\circ$}
\end{axis}
\end{tikzpicture}

\caption{$C_L(\alpha)$ non visqueux de trois profils (traits pleins) et théorie des profils minces (tirets), avec l'angle $\alpha_{L0}$ calculé par l'intégrale du chapitre~\ref{ch:potentiel}. L'écart de pente vient de l'épaisseur.}\label{fig:cl-alpha}
\end{figure}

\begin{resultat}
\begin{itemize}
\item $N$ sources constantes, un tourbillon uniforme, $N$ conditions de glissement et une condition de Kutta : un système dense $(N+1)\times(N+1)$, inversé une fois par géométrie.
\item Sorties : $V_t$ signée, $C_p$, $C_L$, $C_M$, et deux contrôles ($C_{L,KJ}$, $C_{D,p}$).
\item Précision : $1/N$, soit moins de 0,1~\% sur $C_L$ à 200 panneaux.
\end{itemize}
\end{resultat}

\begin{projet}
Côté optimisation, deux propriétés de ce bloc comptent. D'abord, il est \emph{exactement} linéaire en $\alpha$ pour une géométrie donnée, ce qui rend $C_L$ quasi affine en $(m,\alpha)$ : un modèle quadratique représente $C_L$ avec $R^2=0{,}9998$. Ensuite, c'est une méthode d'éléments de frontière, comme la méthode des moments qu'utilisent beaucoup de solveurs de SER : on discrétise la surface et on résout un système dense. C'est aussi le type de système que visent les algorithmes quantiques de résolution linéaire cités dans la note d'intention.
\end{projet}

\exercices
\begin{exercice}
Vérifier le calcul de $v_\sigma$ : montrer que $\int_0^L\frac{\eta\,\dd\xi'}{(\xi-\xi')^2+\eta^2}=\theta_2-\theta_1$, avec $\theta_k=\operatorname{atan2}(\eta,\xi-\xi_k)$. Que vaut cette intégrale très loin du panneau ? Interpréter.
\end{exercice}
\begin{exercice}
(a) Quel est le coût, en ordre de grandeur, de l'inversion d'une matrice $201\times201$, puis d'une résolution pour une nouvelle incidence une fois l'inverse connue ? (b) En déduire pourquoi \code{alpha\_for\_cl}, qui appelle \code{solve} 5 à 8 fois, ne coûte presque rien de plus qu'un seul calcul.
\end{exercice}
